\documentclass[11pt,a4paper]{report}

% ============================================================
%  IBM "RAG and Agentic AI" Spezialisierung — Kursdokumentation
%  Ein Kapitel pro Kurs (10 Kurse insgesamt).
% ============================================================

\usepackage[ngerman]{babel}
\usepackage[utf8]{inputenc}
\usepackage[T1]{fontenc}
\usepackage[a4paper,margin=2.5cm]{geometry}
\usepackage{hyperref}
\usepackage{xcolor}
\usepackage{listings}
\usepackage{enumitem}
\usepackage{tcolorbox}
\usepackage{booktabs}
\usepackage{graphicx}
\usepackage{titlesec}
\usepackage{parskip}
\usepackage[strict=true]{csquotes}
\usepackage{seqsplit}
\usepackage[htt]{hyphenat}

\tcbuselibrary{skins,breakable}

% --- Code-Stil (Python, LangChain-Snippets etc.) ---
\lstdefinestyle{code}{
    basicstyle=\ttfamily\small,
    keywordstyle=\color{blue!70!black}\bfseries,
    commentstyle=\color{gray},
    stringstyle=\color{teal!70!black},
    breaklines=true,
    breakatwhitespace=false,
    numbers=left,
    numberstyle=\tiny\color{gray},
    frame=single,
    rulecolor=\color{gray!50},
    backgroundcolor=\color{gray!5},
    tabsize=4,
    showstringspaces=false
}
\lstset{
    style=code,
    literate=%
        {Ä}{{\"A}}1 {Ö}{{\"O}}1 {Ü}{{\"U}}1
        {ä}{{\"a}}1 {ö}{{\"o}}1 {ü}{{\"u}}1
        {ß}{{\ss}}1
}

% --- Box "Das Wichtigste in Kürze" (wird am Ende jedes Kapitels benutzt) ---
\newtcolorbox{merke}[1][]{%
    colback=blue!5,
    colframe=blue!40!black,
    fonttitle=\bfseries,
    title=Das Wichtigste in Kürze,
    breakable,
    #1
}

% --- Box für offene Punkte / Lücken im Quellmaterial ---
\newtcolorbox{luecke}[1][]{%
    colback=orange!5,
    colframe=orange!60!black,
    fonttitle=\bfseries,
    title=Anmerkung zum Quellmaterial,
    breakable,
    #1
}

\hypersetup{
    colorlinks=true,
    linkcolor=blue!50!black,
    urlcolor=blue!50!black,
    citecolor=blue!50!black,
    pdftitle={Kursdokumentation -- IBM RAG and Agentic AI Spezialisierung},
    pdfauthor={Zidane K.}
}

\title{Kursdokumentation\\[0.3em]
       \large IBM ``RAG and Agentic AI'' Spezialisierung (Coursera)}
\author{Zidane K.}
\date{\today}

\begin{document}

\maketitle
\tableofcontents

\chapter{Kurs 1 --- Generative-KI-Anwendungen entwickeln: Erste Schritte}

\section{Überblick und Lernziele}

Dieser erste Kurs der Spezialisierung legt das Fundament für alles, was danach kommt:
Retrieval-Augmented Generation (RAG), Vektordatenbanken und agentenbasierte Systeme bauen
alle auf denselben Grundbausteinen auf, die hier eingeführt werden. Der Kurs beantwortet
im Kern drei Fragen: Wie spricht man ein Foundation Model überhaupt sinnvoll an (Prompt
Engineering)? Wie strukturiert und automatisiert man diese Ansprache, statt jedes Mal von
Hand zu formulieren (LangChain-Bausteine wie \texttt{PromptTemplate}, Chains, Parser)? Und
wie baut man daraus eine erste, echte Anwendung, die mehr kann als nur Text zurückgeben
(eine Flask-Webanwendung mit strukturierter, verlässlicher Ausgabe)?

Am Ende des Kurses soll man:
\begin{itemize}
    \item verstehen, wo ein LLM in einer Anwendungsarchitektur sitzt und wie man es über
          eine API/ein SDK anspricht (Prompt, Parameter, Antwort);
    \item mehrere Prompt-Engineering-Techniken (Zero-/One-/Few-Shot, Chain-of-Thought,
          Self-Consistency) kennen und wissen, wann welche sinnvoll ist;
    \item mit LangChain Prompts, Parser, Dokumente, Vektorspeicher, Retriever, Gedächtnis,
          Chains und Agenten zu einer funktionierenden Pipeline zusammensetzen können;
    \item eine erste eigenständige generative KI-Anwendung (Flask-Webapp) bauen können, die
          strukturierte statt roher Textantworten liefert.
\end{itemize}

Alle praktischen Teile des Kurses wurden gegen IBM-watsonx.ai-Modelle umgesetzt
(u.\,a. \texttt{\seqsplit{ibm/granite-4-h-small}}, \texttt{\seqsplit{meta-llama/llama-4-maverick-17b-128e-instruct-fp8}},
\texttt{\seqsplit{mistralai/mistral-small-3-1-24b-instruct-2503}}), angesprochen
über die Bibliotheken \texttt{ibm-watsonx-ai} und \texttt{langchain-ibm}.

\section{Kernkonzepte}

\subsection{Foundation Models und Interaktion über watsonx.ai}

Ein Foundation Model ist die Kernkomponente jeder generativen KI-Anwendung: Es nimmt
Klartext als Eingabe und erzeugt Klartext als Ausgabe. Alles, was eine Anwendung
\enquote{intelligent} macht, entsteht dadurch, wie man dieses Modell einbettet und ansteuert ---
nicht durch das Modell allein. In der Praxis bedeutet das: Man instanziiert ein Modell
über \texttt{ModelInference} (aus \texttt{ibm-watsonx-ai}) mit drei Angaben --- einer
\texttt{model\_id} (welches Modell), \texttt{parameters} (wie das Modell generiert) und
\texttt{credentials}/\texttt{project\_id} (Zugang) --- und ruft es dann mit
\texttt{model.generate(prompt)} auf. Die wichtigsten Generierungsparameter sind
\texttt{MAX\_NEW\_TOKENS}/\texttt{MIN\_NEW\_TOKENS} (Länge der Antwort) und
\texttt{TEMPERATURE}/\texttt{TOP\_P}/\texttt{TOP\_K} (Zufälligkeit bzw. Kreativität der
Antwort). Eine niedrige Temperatur (z.\,B. 0{,}1--0{,}2) liefert deterministische,
konsistente Antworten und eignet sich für faktenbasierte oder anweisungsbefolgende
Aufgaben; eine hohe Temperatur (z.\,B. 0{,}8) liefert vielfältigere, kreativere Antworten,
büßt dafür an Konsistenz zwischen mehreren Durchläufen ein.

Um ein watsonx.ai-Modell in LangChain nutzbar zu machen (Chains, Prompt-Templates, etc.),
wird es mit \texttt{WatsonxLLM} umschlossen (\texttt{WatsonxLLM(model=model)}). Danach
verhält sich das Modell wie jedes andere LangChain-LLM und kann in Chains eingebunden
werden.

\subsection{Chat-Modelle und Nachrichtentypen}

Ein Chat-Modell unterscheidet sich von einem einfachen Text-Modell dadurch, dass es
Rollen kennt: eine Konversation besteht aus einer Liste von Nachrichten, die jeweils eine
Rolle und einen Inhalt haben. Die drei wichtigsten Typen sind \texttt{SystemMessage}
(setzt das Verhalten/die Persona des Modells, meist die erste Nachricht),
\texttt{HumanMessage} (die Eingabe eines Nutzers) und \texttt{AIMessage} (eine Antwort
des Modells). Reicht man dem Modell eine Liste solcher Nachrichten, kann es auf den
gesamten bisherigen Gesprächsverlauf reagieren --- das ist die Grundlage jeder
Chat-Anwendung mit Kontext über mehrere Runden.

\subsection{Prompting-Strategien: Zero-, One- und Few-Shot}

Diese drei Begriffe beschreiben, wie viele Beispiele man dem Modell im Prompt mitgibt,
bevor man die eigentliche Frage stellt:
\begin{itemize}
    \item \textbf{Zero-Shot}: Es wird kein Beispiel mitgegeben, nur die Aufgabe selbst
          (\enquote{Fasse folgenden Text zusammen: ...}). Funktioniert für einfache, allgemein
          bekannte Aufgaben gut, ist aber anfällig für ein unerwartetes Antwortformat.
    \item \textbf{One-Shot}: Ein einziges Beispiel wird vorangestellt, um Format und Ton
          der gewünschten Antwort zu demonstrieren.
    \item \textbf{Few-Shot}: Mehrere Beispiele werden vorangestellt. Das stabilisiert das
          Antwortformat deutlich, kostet aber Prompt-Länge (und damit Tokens/Kosten).
\end{itemize}
Die Praxisregel lautet: so viele Beispiele wie nötig, so wenige wie möglich --- Beispiele
werden nur so lange hinzugefügt, bis sich das Format des Modells stabilisiert.

\subsection{Chain-of-Thought und Self-Consistency}

Bei mehrstufigen Denkaufgaben (z.\,B. Rechenaufgaben, logische Schlussfolgerungen) hilft
es, das Modell explizit aufzufordern, seinen Gedankengang Schritt für Schritt
offenzulegen, statt direkt eine Endantwort zu liefern (\emph{Chain-of-Thought},
z.\,B. durch einen Zusatz wie \enquote{Denke Schritt für Schritt}). Das verbessert die
Genauigkeit bei komplexeren Aufgaben spürbar, weil das Modell seine eigene Herleitung
\enquote{sieht} und darauf aufbauen kann, statt in einem Schritt zu raten.

\emph{Self-Consistency} baut darauf auf: Man lässt das Modell dieselbe Chain-of-Thought-
Aufgabe mehrfach mit unterschiedlicher Zufälligkeit (höhere Temperatur) durchlaufen und
nimmt am Ende die häufigste (mehrheitliche) Antwort über alle Durchläufe. Das reduziert
das Risiko, dass ein einzelner fehlerhafter Gedankengang die Endantwort verfälscht.

\subsection{LangChain PromptTemplates}

Ein \texttt{PromptTemplate} trennt den festen Wortlaut eines Prompts von den variablen
Daten, die eingesetzt werden. Statt für jede Eingabe einen neuen String
zusammenzubauen, definiert man eine Vorlage mit Platzhaltern
(\lstinline|PromptTemplate.from_template("Erzähl mir einen {adjective} Witz über {topic}")|)
und füllt sie bei Bedarf mit einem Dictionary
(\lstinline|{"adjective": "lustigen", "topic": "Katzen"}|). Das macht Prompts
wiederverwendbar und testbar --- ein Template bedient viele Eingaben, statt für jede neue
Eingabe von Hand einen Prompt zu schreiben. \texttt{ChatPromptTemplate} tut dasselbe für
ganze Nachrichtenlisten (System-/Human-Nachrichten mit Platzhaltern), und
\texttt{MessagesPlaceholder} erlaubt es, an einer bestimmten Stelle im Template eine
ganze, zur Laufzeit übergebene Liste von Nachrichten einzufügen (z.\,B. den bisherigen
Gesprächsverlauf).

\subsection{Output Parser: von Freitext zu strukturierten Daten}

Ein LLM liefert von sich aus Freitext zurück. Für viele Anwendungen (z.\,B. eine
Weboberfläche, die bestimmte Felder anzeigen will) reicht das nicht --- man braucht
strukturierte, verlässlich geformte Daten. Ein Output Parser übernimmt genau diese
Übersetzung. Die wichtigsten Varianten:
\begin{itemize}
    \item \texttt{JsonOutputParser} (ggf. mit einem Pydantic-\texttt{BaseModel}): erzwingt
          ein bestimmtes JSON-Schema. Der Parser erzeugt dazu automatisch
          Formatierungsanweisungen (\texttt{get\_format\_instructions()}), die man in den
          Prompt einbettet, damit das Modell weiß, welches Format erwartet wird.
    \item \texttt{CommaSeparatedListOutputParser}: parst eine kommagetrennte Antwort in
          eine Python-Liste.
    \item \texttt{StrOutputParser}: extrahiert einfach den reinen Text-String aus der
          Modellantwort, nützlich als letztes Glied einer LCEL-Chain.
\end{itemize}
Diese Kombination aus Prompt-Formatinstruktionen und einem passenden Parser ist der
Mechanismus, mit dem ein LLM zu einer zuverlässigen Datenquelle für den Rest einer
Anwendung wird.

\subsection{Retrieval-Augmented Generation: Dokumente, Splitter, Embeddings, Vektorspeicher, Retriever}

RAG löst ein grundlegendes Problem von LLMs: Ihr Wissen ist auf die Trainingsdaten
begrenzt und sie kennen private/aktuelle Dokumente nicht. RAG-Pipelines bestehen aus
einer festen Abfolge von Bausteinen:
\begin{enumerate}
    \item \textbf{Document Loader} lädt Rohdaten (z.\,B. \texttt{PyPDFLoader} für PDFs,
          \texttt{WebBaseLoader} für Webseiten) in \texttt{Document}-Objekte
          (\texttt{page\_content} + \texttt{metadata}).
    \item \textbf{Text Splitter} zerlegt lange Dokumente in kleinere Chunks, die in das
          Kontextfenster des Modells passen (z.\,B. \texttt{CharacterTextSplitter},
          \texttt{RecursiveCharacterTextSplitter}), mit einstellbarer Chunk-Größe und
          Überlappung zwischen Chunks (damit an Chunk-Grenzen kein Kontext verloren geht).
    \item \textbf{Embedding-Modell} (z.\,B. \texttt{WatsonxEmbeddings} mit
          \texttt{\seqsplit{ibm/granite-embedding-278m-multilingual}}) wandelt jeden Chunk in einen
          Vektor um, der die Bedeutung des Texts in einem hochdimensionalen Raum
          repräsentiert.
    \item \textbf{Vektorspeicher} (z.\,B. \texttt{Chroma}) speichert diese Vektoren und
          erlaubt eine Ähnlichkeitssuche: Man embeddet eine Suchanfrage genauso und
          erhält die Chunks, deren Vektoren am nächsten liegen.
    \item \textbf{Retriever} ist die allgemeine Schnittstelle darüber (z.\,B.
          \texttt{vectorstore.as\_retriever()}) --- er nimmt eine Text-Query entgegen und
          liefert passende \texttt{Document}-Objekte zurück, unabhängig davon, wie sie
          gespeichert sind. Ein \texttt{ParentDocumentRetriever} löst dabei einen
          Zielkonflikt: kleine Chunks liefern präzisere Embeddings, größere Chunks
          erhalten mehr Kontext --- dieser Retriever sucht über kleine Chunks, gibt aber
          die zugehörigen größeren \enquote{Eltern-Chunks} zurück.
    \item \textbf{RetrievalQA} verbindet Retriever und LLM zu einer fertigen
          Frage-Antwort-Kette: Die Retriever-Treffer werden automatisch in den Prompt an
          das LLM eingefügt (\texttt{stuff}), das daraus die Antwort formuliert.
\end{enumerate}
Kurz: RAG ist load $\rightarrow$ split $\rightarrow$ embed $\rightarrow$ store
$\rightarrow$ retrieve $\rightarrow$ in den Prompt einfügen $\rightarrow$ generieren.

\subsection{Gedächtnis (Memory) in Konversationen}

Ein Chat-Modell selbst hat kein Gedächtnis zwischen zwei Aufrufen --- jeder Aufruf ist
zustandslos. Damit eine Anwendung sich an frühere Nachrichten \enquote{erinnert}, muss der
bisherige Verlauf bei jedem neuen Aufruf explizit mitgegeben werden. \texttt{
ChatMessageHistory} ist dafür die einfachste Bausteinklasse: Sie sammelt
\texttt{HumanMessage}s und \texttt{AIMessage}s und kann diese Liste dem Modell erneut
übergeben. \texttt{ConversationBufferMemory} wickelt das automatisch in eine
\texttt{ConversationChain} ein: Bei jedem neuen Nutzer-Input wird automatisch der
gesamte bisherige Verlauf in den Prompt eingefügt.
\texttt{ConversationSummaryMemory} macht dasselbe, komprimiert den Verlauf aber
fortlaufend zu einer Zusammenfassung (per eigenem LLM-Aufruf) statt ihn vollständig zu
wiederholen --- das hält den Prompt bei langen Gesprächen kürzer, auf Kosten von Detail.

\subsection{Chains: klassisch vs. LCEL}

Eine Chain verkettet mehrere Schritte (Prompt $\rightarrow$ LLM $\rightarrow$ Parser
$\rightarrow$ ggf. nächster Prompt) zu einem Ablauf, der mit einem Aufruf ausgeführt
wird. Es gibt zwei Bauweisen, die im Kurs bewusst nebeneinander gezeigt wurden:
\begin{itemize}
    \item \textbf{Klassisch} (\texttt{LLMChain}, \texttt{SequentialChain}): jeder
          Schritt ist ein explizites Objekt mit einem benannten \texttt{output\_key};
          mehrere Chains werden über diese Namen zu einer \texttt{SequentialChain}
          zusammengesteckt. Funktional, aber vergleichsweise starr.
    \item \textbf{LCEL} (LangChain Expression Language): Schritte werden mit dem
          Pipe-Operator verbunden, z.\,B. \texttt{chain = prompt | llm | parser}.
          Mehrstufige Abläufe, bei denen der Output eines Schritts in den nächsten
          einfließt, werden mit \texttt{RunnablePassthrough.assign(...)} gebaut: Man übergibt
          ihm einen neuen Schlüsselnamen und eine \texttt{lambda}-Funktion, die aus dem
          bisherigen Eingabe-Dictionary den nächsten Wert berechnet. So werden dem
          Dictionary schrittweise neue Felder hinzugefügt, ohne bestehende zu überschreiben.
\end{itemize}
LCEL gilt inzwischen als Standardweg für neue Chains; \texttt{SequentialChain} wird nur
noch aus Kompatibilitätsgründen unterstützt.

\subsection{Tools und Agenten (ReAct)}

Ein LLM kann von sich aus nur Text erzeugen --- es kann nicht rechnen, nicht im Internet
suchen, nicht auf externe Systeme zugreifen. Ein \textbf{Tool} ist eine Funktion mit
Namen und Beschreibung (z.\,B. ein Python-Rechner über \texttt{PythonREPL}, oder eine
eigene Funktion via \texttt{@tool}-Dekorator), die ein Agent bei Bedarf aufrufen kann.
Ein \textbf{Agent} nutzt das LLM als Entscheidungsinstanz: Er überlegt, welches Tool zu
einer Anfrage passt, ruft es mit passender Eingabe auf, liest das Ergebnis und
entscheidet, ob noch ein weiterer Schritt nötig ist oder die Endantwort feststeht. Das
im Kurs verwendete Muster ist \textbf{ReAct} (Reasoning + Acting): Der Agent durchläuft
wiederholt den Zyklus \emph{Thought} (was tun?) $\rightarrow$ \emph{Action} (welches
Tool, welche Eingabe?) $\rightarrow$ \emph{Observation} (Ergebnis des Tools), bis er
mit \emph{Final Answer} abschließt. Technisch wird das mit
\texttt{create\_react\_agent(llm, tools, prompt)} gebaut und von einem
\texttt{AgentExecutor} ausgeführt, der diese Schleife orchestriert und Fehler beim
Parsen der Modellantwort abfangen kann (\texttt{handle\_parsing\_errors=True}).

\section{Praktische Umsetzung}

\subsection{Lab 1 --- Prompt Engineering und LangChain PromptTemplates}

In diesem Lab wurde \texttt{ibm/granite-4-h-small} über \texttt{langchain-ibm}s
\texttt{WatsonxLLM} in einer wiederverwendbaren Hilfsfunktion
\texttt{llm\_model(prompt, params)} gekapselt. Anhand identischer Aufgaben wurden
Zero-Shot, One-Shot und Few-Shot verglichen, danach Chain-of-Thought für mehrstufige
Denkaufgaben eingesetzt und mit Self-Consistency (mehrere Durchläufe, Mehrheitsantwort)
kombiniert. Abschließend wurden hart kodierte Prompts durch \texttt{PromptTemplate}s mit
Eingabevariablen ersetzt und mit dem Modell zu einer Chain verbunden. Fünf Übungen
vertieften das an Aufgaben wie Zusammenfassung, Frage-Antwort, Klassifikation,
Code-Generierung und Rollenspiel.

\subsection{Lab 2 --- LangChain-Grundbausteine (RAG, Memory, Chains, Agenten)}

Dieses Lab deckte in sieben Übungen den kompletten LangChain-Werkzeugkasten ab:

\begin{itemize}
    \item Modelle und Chat-Nachrichten; String- und Chat-\texttt{PromptTemplate}s samt
          \texttt{MessagesPlaceholder}.
    \item \texttt{JsonOutputParser} und \texttt{CommaSeparatedListOutputParser}.
    \item Eine vollständige RAG-Pipeline: PDF-/Web-Loader, zwei Text-Splitter im
          Vergleich, \texttt{WatsonxEmbeddings}, \texttt{Chroma} als Vektorspeicher,
          ein Vektorspeicher- sowie ein \texttt{ParentDocumentRetriever}, und
          \texttt{RetrievalQA} als fertige Frage-Antwort-Kette.
    \item Gesprächsgedächtnis: \texttt{ChatMessageHistory} sowie
          \texttt{ConversationBufferMemory} und \texttt{ConversationSummaryMemory}.
    \item Chains klassisch (\texttt{LLMChain}/\texttt{SequentialChain}) vs. LCEL.
    \item Ein ReAct-Agent mit einem Python-Rechner-Tool und einem eigenen
          Text-Formatierungs-Tool.
\end{itemize}

\subsection{Abschlussprojekt --- GenAI Flask App}

Das Abschlussprojekt (\enquote{Build Your First GenAI Application The Right Way}) bündelt die
Kurskonzepte zu einer eigenständigen Flask-Webanwendung (\enquote{AI Assistant}): Ein Nutzer
wählt im Frontend ein Modell aus (Granite, Llama oder Mistral) und stellt eine Frage; der
Flask-Endpunkt \texttt{/generate} leitet sie an das gewählte Modell weiter. Die zentrale
Neuerung gegenüber den Labs ist, dass die Antwort nicht mehr Freitext ist, sondern über
ein Pydantic-Modell \texttt{AIResponse} (\texttt{summary}, \texttt{sentiment},
\texttt{response}, \texttt{category}, \texttt{action}) und einen
\texttt{JsonOutputParser} als validiertes JSON erzwungen wird --- geeignet für einen
Kundensupport-Anwendungsfall, in dem eine nachgelagerte Anwendung mit diesen Feldern
weiterarbeiten könnte. Da jedes Modell eine andere native Prompt-Syntax erwartet
(Llamas Spezial-Tokens \texttt{<|begin\_of\_text|>}/\texttt{<|start\_header\_id|>},
ein einfaches \texttt{System:}/\texttt{Human:}-Layout für Granite, Mistrals
\texttt{[INST]}-Tags), bekommt jedes Modell sein eigenes \texttt{PromptTemplate} bei
identischen Eingabevariablen. Eine Übung erweiterte das Schema zusätzlich um
\texttt{category} und \texttt{action}, um dem Support-Mitarbeiter eine konkrete
Handlungsempfehlung mitzugeben.

\section{Wichtige Zusammenhänge und Fallstricke}

\begin{itemize}
    \item Prompt-Struktur hat oft einen größeren Effekt auf die Antwortqualität als
          Parameter-Tuning: Bevor an \texttt{temperature}/\texttt{top\_p}/\texttt{top\_k}
          gedreht wird, lohnt es sich meist mehr, das Prompt-Format (Beispiele, klare
          Formatanweisungen) zu verbessern.
    \item \texttt{JsonOutputParser} funktioniert nur zuverlässig, wenn die
          Formatinstruktionen (\texttt{get\_format\_instructions()}) tatsächlich in den
          Prompt eingebettet werden --- ein Modell \enquote{errät} das gewünschte JSON-Schema
          sonst nicht zuverlässig.
    \item Ein Prompt-Template, das für ein Modell funktioniert, überträgt sich nicht
          automatisch auf ein anderes: Jede Modellfamilie hat ihr eigenes natives
          Nachrichtenformat (siehe Abschlussprojekt); das im Kurs mehrfach gezeigte
          Muster ist, dieselben logischen Eingaben (\texttt{system\_prompt},
          \texttt{user\_prompt}) in mehrere modellspezifische Templates zu gießen.
    \item RAG und Agenten setzen beide auf denselben Grundbausteinen auf (Prompt, LLM,
          Parser) --- RAG ist eine feste Pipeline ohne Entscheidungslogik, ein Agent
          entscheidet dagegen zur Laufzeit, welchen Schritt er als Nächstes braucht.
    \item \texttt{SequentialChain} und LCEL lösen dieselbe Aufgabe; LCEL wird für neue
          Chains empfohlen, \texttt{SequentialChain} taucht in älterem Code/Beispielen
          aber weiterhin auf und sollte lesbar bleiben.
\end{itemize}

\begin{luecke}
Die genauen Coursera-Modulnamen und die exakte Gliederung der Video-Lektionen lagen für
dieses Kapitel nicht vollständig vor; die Struktur oben folgt der Reihenfolge, in der die
Konzepte in den Labs und dem Abschlussprojekt tatsächlich verwendet wurden.
\end{luecke}

\section{Zusammenfassung: Das Wichtigste in Kürze}

\begin{merke}
\begin{itemize}[leftmargin=1.2em]
    \item Ein Foundation Model wird über \texttt{model\_id}, \texttt{parameters}
          (v.\,a. Temperatur) und \texttt{credentials} angesprochen; für LangChain wird es
          mit \texttt{WatsonxLLM} umschlossen.
    \item Few-Shot-Beispiele stabilisieren das Antwortformat; Chain-of-Thought plus
          Self-Consistency verbessert mehrstufige Denkaufgaben.
    \item \texttt{PromptTemplate} trennt Prompt-Wortlaut von variablen Daten und macht
          Prompts wiederverwendbar; \texttt{ChatPromptTemplate}/\texttt{
          MessagesPlaceholder} tun dasselbe für Nachrichtenlisten.
    \item Ein Output Parser (v.\,a. \texttt{JsonOutputParser} mit einem Pydantic-Schema)
          verwandelt Freitext in strukturierte, verlässliche Daten --- die Grundlage
          jeder echten Anwendung, die mit der LLM-Antwort weiterarbeiten will.
    \item RAG ist die feste Pipeline load $\rightarrow$ split $\rightarrow$ embed
          $\rightarrow$ store $\rightarrow$ retrieve $\rightarrow$ generate;
          \texttt{RetrievalQA} kapselt das Ganze in eine Zeile Code.
    \item Gedächtnis muss aktiv verwaltet werden (\texttt{ChatMessageHistory},
          \texttt{ConversationBufferMemory}/\texttt{ConversationSummaryMemory}) --- ein
          LLM-Aufruf ist für sich genommen zustandslos.
    \item LCEL (\texttt{prompt | llm | parser}) ist der moderne, empfohlene Weg, Chains
          zu bauen; \texttt{SequentialChain} bleibt aus Kompatibilitätsgründen bestehen.
    \item Ein Agent (z.\,B. ReAct) unterscheidet sich von einer Chain dadurch, dass er
          zur Laufzeit entscheidet, welches Tool er wie oft aufruft, statt einen fest
          vorgegebenen Ablauf abzuarbeiten.
\end{itemize}
\end{merke}

\chapter{Kurs 2 --- Build RAG Applications: Get Started}

\begin{luecke}
Kursinhalt noch nicht bereitgestellt. Sobald die Informationen zu diesem Kurs vorliegen,
wird dieses Kapitel gemäß dem Prompt in \texttt{PROMPT.md} ausgefüllt.
\end{luecke}

% \section{Überblick und Lernziele}
% \section{Kernkonzepte}
% \section{Praktische Umsetzung}
% \section{Wichtige Zusammenhänge und Fallstricke}
% \section{Zusammenfassung: Das Wichtigste in Kürze}

\input{chapters/kapitel03_vector_databases}
\input{chapters/kapitel04_advanced_rag_retrievers}
\chapter{Kurs 5 --- Build Multimodal Generative AI Applications}

\begin{luecke}
Kursinhalt noch nicht bereitgestellt. Sobald die Informationen zu diesem Kurs vorliegen,
wird dieses Kapitel gem\"a\ss{} dem Prompt in \texttt{PROMPT.md} ausgef\"ullt.
\end{luecke}

% \section{\"Uberblick und Lernziele}
% \section{Kernkonzepte}
% \section{Praktische Umsetzung}
% \section{Wichtige Zusammenh\"ange und Fallstricke}
% \section{Zusammenfassung: Das Wichtigste in K\"urze}

\input{chapters/kapitel06_fundamentals_ai_agents}
\input{chapters/kapitel07_agentic_langchain_langgraph}
\chapter{Kurs 8 --- Agentic AI with LangGraph, CrewAI, AutoGen and BeeAI}

\begin{luecke}
Kursinhalt noch nicht bereitgestellt. Sobald die Informationen zu diesem Kurs vorliegen,
wird dieses Kapitel gem\"a\ss{} dem Prompt in \texttt{PROMPT.md} ausgef\"ullt.
\end{luecke}

% \section{\"Uberblick und Lernziele}
% \section{Kernkonzepte}
% \section{Praktische Umsetzung}
% \section{Wichtige Zusammenh\"ange und Fallstricke}
% \section{Zusammenfassung: Das Wichtigste in K\"urze}

\input{chapters/kapitel09_ai_agents_mcp}
\input{chapters/kapitel10_capstone_project}

\end{document}
